%% kapa: exam=1 round=27 year=2016 period=1 form=A subject=real_estate subject_ko=부동산학원론 seq=38 number=118
%% source: Q-Net 감정평가사 기출문제 (cst00309) artlSeq=5169402 file=extracted/1008/5169402/제27회 감정평가사 1차 문제/제27회 감정평가사 1교시문제(A형).pdf
%% conversion: kordoc 4.24.51 -> tools/gen1.py (auto); review=proofread-sonnet; flags=none
\begin{question}{118}
부동산관리에 관한 설명으로 옳지 않은 것은?
\begin{choices}
  \item 자산관리(Asset Management)는 부동산자산을 포트폴리오(portfolio) 관점에서 관리하는 자산․부채의 종합관리를 의미한다.
  \item 재산관리(Property Management)는 시설사용자나 사용과 관련한 타부문의 요구에 단순히 부응하는 정도의 소극적이고 기술적인 측면을 중시하는 부동산관리를 의미한다.
  \item 대상건물의 기능을 유지하기 위해서 건물에 대해 수리 및 점검을 하는 등의 관리는 기술적 측면의 관리에 해당한다.
  \item 위탁관리방식은 전문업자를 이용함으로써 합리적이고 편리하며, 전문화된 관리와 서비스를 받을 수 있다는 장점이 있다.
  \item 기밀유지 측면에서는 자가관리방식이 위탁관리방식보다 유리하다.
\end{choices}
\end{question}
\answer{2}
